\chapter{Governance, Risks, and Open Decisions}
\ChapterLabel{DevelopmentPlan}
\label{Section:DevelopmentPlanResponsibility}
This chapter defines project ownership, decision gates, tailoring, and change control.


\section{Project Authorship}
\label{Section:ProjectAuthorship}
Project documentation is co-authored by \DocumentAuthors.
The user confirmed Chadwick Adebiyi as co-author of all project documents
on November 23, 2026. The authorship metadata applies to this integrated
specification and to project-authored supporting documents. Ownership,
decision coordination, and acceptance authority are recorded separately in
\hyperref[Table:ProjectOwners]{the owner register}.

\section{Confirmed Governance}
The user reconfirmed these assignments for SPE-800-Sensor on November 9, 2026. These roles remain in effect until changed by the project owner. Recording these assignments does not claim that technical acceptance has occurred.
\small
\begin{tabularx}{\linewidth}{L{4.0cm}Y}
\caption{Project Owners\label{Table:ProjectOwners}}\\
\toprule
Role & Assigned owner\\\midrule
\endfirsthead
\toprule
Role & Assigned owner\\\midrule
\endhead
\bottomrule
\endlastfoot
Document owner & Michael Stoddard\\
Organization & Stevens Institute\\
Technical owner & Michael Stoddard\\
Verification/validation owner & Michael Stoddard\\
Acceptance authority & Dr. Darian Muresan\\
\end{tabularx}
\normalsize

\Needspace{12\baselineskip}
\section{Domain Profile Selection}
\label{Section:DevelopmentPlanProfiles}
\small
\begin{tabularx}{\linewidth}{L{3cm}L{2.4cm}Y}
\caption{Project Domain Profile Decisions\label{Table:ProjectDomainProfiles}}\\
\toprule
Profile & Status & Rationale and decision provenance\\\midrule
\endfirsthead
\toprule
Profile & Status & Rationale and decision provenance\\\midrule
\endhead
\bottomrule
\endlastfoot
Embedded Systems & Applicable & Confirmed by the user November 9, 2026: the system includes sensor hardware, conversion, timing, and a hardware/software interface. Technical tailoring owner: Michael Stoddard. This is the only supplied profile; additional profiles require a separate decision.\\
\end{tabularx}
\normalsize

\section{Stage and Change Control}
This stage covers documentation only. The user confirmed the namespace \texttt{s8s}. Authored sources and editable diagrams in \path|Architecture/| are authoritative for this project; the reusable template remains separate. Independent technical review and acceptance remain human work. Repository authoring and AI execution guidance are maintained outside this publication.

Maintain the architecture in \path|Architecture/| within the SPE-800 repository. Changes use a branch and pull request against \texttt{main}. This candidate is pending document review. Michael Stoddard coordinates the document review; Dr. Darian Muresan retains technical acceptance authority. No approval or publication is inferred from a successful build.

\section{Owned Questions and Decision Gates}
\label{Section:OwnedQuestions}
Question ownership follows the supplied roles; it assigns responsibility for obtaining an answer and does not claim the owner has answered it. Each question below remains Open for the residual decisions explicitly stated. Supplied facts are acknowledged in the question rather than silently treating the entire topic as unresolved. Dependent design and implementation wait for their disposition. Decisions belong here once, and other chapters link to them.
\small
\begin{tabularx}{\linewidth}{L{1.3cm}L{3.1cm}Y}
\caption{Owned Questions\label{Table:OwnedQuestions}}\\
\toprule
ID & Owner & Question and impact\\\midrule
\endfirsthead
\toprule
ID & Owner & Question and impact\\\midrule
\endhead
\bottomrule
\endlastfoot
\phantomsection\label{Question:Q-101}Q-101 & Michael Stoddard & Confirmed: mission is measurement usability; stakeholders are Operator, Acquisition-Developer, Field-Engineer, and Mission-Integrator. Define the intended operating context and measurable stakeholder success for SPE-800-Sensor. Blocks final validation criteria.\\
\phantomsection\label{Question:Q-102}Q-102 & Michael Stoddard & Confirmed: ADC0804 produces an 8-bit conversion code; FT245R supplies the USB-to-parallel path. Select sensor/physical quantity, signal range and units, ADC reference and conversion clock/handshake, calibration, exact device/board revisions, and electrical compatibility. Blocks the schematic and analog acceptance limits.\\
\phantomsection\label{Question:Q-103}Q-103 & Michael Stoddard & Confirmed: two logical parallel inputs feed the interface logic and FT245R USB data path. Define bus ownership, strobes/handshake, transfer order, framing, record boundaries, error detection, loss and backpressure behavior. Blocks a realizable transfer contract; see \hyperref[Question:Q-112]{Q-112}.\\
\phantomsection\label{Question:Q-104}Q-104 & Michael Stoddard & Set sampling rate, capture skew, end-to-end latency, jitter, queue/resource budgets, power limits, and operating conditions with units and tolerances. Hardware selection alone does not supply project performance limits.\\
\phantomsection\label{Question:Q-105}Q-105 & Michael Stoddard & Confirmed: the host qualifies measurements and publishes code, status, and timing. Define the human interaction form, Operator decision presentation, Field-Engineer access, and whether configuration, storage, or reverse commands are needed. No GUI or command channel is selected.\\
\phantomsection\label{Question:Q-106}Q-106 & Michael Stoddard & Confirmed: the current host is macOS; Windows/Linux are outside this revision. Select macOS versions, CPU architectures, C++ standard, compilers, libraries, FT245R driver/API, permissions, packaging, and build tools. Blocks reproducible host builds.\\
\phantomsection\label{Question:Q-107}Q-107 & Michael Stoddard & Define startup, shutdown, reset, disconnect, malformed/partial records, invalid data, overflow, stalled operation, degraded behavior, safe states, recovery, and security/privacy boundaries. Fault detection and handling are not implied by the named devices.\\
\phantomsection\label{Question:Q-108}Q-108 & Michael Stoddard & Select reference analog inputs and conversion/counter vectors, fault fixtures, expected qualification/publication results, evidence retention, and representatives of all four stakeholders. Define executable procedures after the interface and decision contracts are accepted.\\
\phantomsection\label{Question:Q-109}Q-109 & Michael Stoddard & Assess manufacturing/provisioning, field updates, maintenance, environment, standards, and component lifecycle for the selected devices; propose coverage or non-applicability dispositions for acceptance-authority approval.\\
\phantomsection\label{Question:Q-110}Q-110 & Michael Stoddard & Propose milestones, risk dispositions, requirement priorities, and the next document baseline version; route this documentation candidate for technical acceptance.\\
\phantomsection\label{Question:Q-111}Q-111 & Michael Stoddard & Define counter width, clock source/frequency/accuracy, capture event, reset and rollover rules, clock-domain crossing, and the relation of counter value to measurement time. Blocks timing interpretation; the counter is not assumed to be 8-bit or an absolute timestamp.\\
\phantomsection\label{Question:Q-112}Q-112 & Michael Stoddard & Define when code and counter are captured, how coherence is ensured, how one association is identified across transfer, and what happens on conversion/counter transitions, reset, partial transfer, or overflow. Compare capture/latch/serialization alternatives after constraints are known; none is selected.\\
\phantomsection\label{Question:Q-113}Q-113 & Michael Stoddard & Define the host qualification rules for measurement usability, all decision inputs, units, freshness/integrity limits, boundary conditions, and unknown/unusable outcomes. Blocks the decision oracle and any usable-measurement claim.\\
\phantomsection\label{Question:Q-114}Q-114 & Michael Stoddard & The Mission application interface uses Open Mission Systems (OMS). Define the sample and sensor-status interfaces for both use cases. The \hyperref[Table:OMSStatusDecisions]{OMS status decisions} govern \RequirementRef{REQ-OMS-01}. Define published code meaning (conversion code versus another code), status semantics, timing fields/origin/units/rollover, association, schema/version/encoding, delivery transport, cadence, and incomplete/error result policy with the downstream consumer. Blocks Mission-Integrator interpretation and publication tests.\\
\phantomsection\label{Question:Q-115}Q-115 & Michael Stoddard & Define what source means (USB device, interface logic, ADC/sensor, clock/counter, or the complete chain), connection evidence, operating-health evidence, diagnostic coverage, reporting mechanism, stale/unknown handling, and false-positive limits. Blocks Field-Engineer acceptance; enumeration and data arrival alone do not establish full source health.\\
\phantomsection\label{Question:Q-116}Q-116 & Michael Stoddard & Allocate interface logic and embedded C++ responsibilities to selected hardware, firmware, or programmable logic, or explicitly revise the earlier embedded-application obligation. Do not infer an MCU, FPGA, or firmware requirement from FT245R selection. Blocks implementation decomposition and embedded-toolchain verification.\\
\phantomsection\label{Question:Q-117}Q-117 & Michael Stoddard & Confirm whether the adjacent \path|mseManual-V1.40.pdf| supplies project evidence or course guidance and identify any relevant sections before incorporating it.\\
\end{tabularx}
\normalsize

\section{OMS Status Interface Decisions}
\RequirementRef{REQ-OMS-01} remains Proposed pending resolution of the
TBD values below. The dates are proposed decision targets. Michael Stoddard
owns coordination of the decisions and records the program-approved values
and their sources before routing the requirement for approval.

\small
\begin{tabularx}{\linewidth}{YL{1.2cm}L{3.5cm}L{3.45cm}}
\caption{OMS Status Interface Decisions\label{Table:OMSStatusDecisions}}\\
\toprule
Decision & Value & Owner & Target decision date\\\midrule
\endfirsthead
\toprule
Decision & Value & Owner & Target decision date\\\midrule
\endhead
\bottomrule
\endlastfoot
OMS standard version & TBD & Michael Stoddard & December 2, 2026\\
CAL profile & TBD & Michael Stoddard & December 2, 2026\\
UCI message & TBD & Michael Stoddard & December 2, 2026\\
Valid, stale, and faulted status mapping & TBD & Michael Stoddard & December 2, 2026\\
Responsible mission gateway & TBD & Michael Stoddard & December 2, 2026\\
Derived sensor-status data product & TBD & Michael Stoddard & December 2, 2026\\
CAL service contract & TBD & Michael Stoddard & December 2, 2026\\
\end{tabularx}
\normalsize

\section{Technical Review Decision Agenda}
\label{Section:ReviewDecisions}
The answers below are TBD. Michael Stoddard owns coordination and records
the agreed answers with the affected stakeholders and acceptance authority.
The OMS status decisions retain their target dates in
\hyperref[Table:OMSStatusDecisions]{the OMS decision register}.

\subsection{High-Priority Decisions}

\paragraph{Q-118. Local and mission sample boundaries.}\label{Question:Q-118}
Who produces and consumes the local record, and who owns mission sample
publication? Is the local consumer the mission gateway, another component,
or an external client? Does the Mission application consume qualified samples
over OMS, and if so, which approved sample contract and delivery endpoint
apply? Which boundary do \RequirementRef{REQ-F-107} and
\RequirementRef{REQ-I-101} govern? Which field meanings, clocks, encodings,
sequence rollover, and invalid or incomplete record dispositions apply at
each boundary? Record the agreed endpoints and contract identifiers.

\paragraph{Q-119. Operator presentation.}\label{Question:Q-119}
Where does the Operator see the reading information? What should appear for
valid, invalid, unavailable, and stale readings, including whether the last
value remains visible? Is age measured from conversion-ready time, sample
capture, or another event? Specify units, display precision, update cadence,
and how unknown age is shown. Which observations and use decisions should
the Operator make in the reference cases for \ValidationRef{VAL-120}?

\paragraph{Q-120. Continuous delivery.}\label{Question:Q-120}
What input rates, burst lengths, run duration, and concurrent processing
loads represent intended operation? What minimum delivery rate, maximum
inter-reading gap, and loss, duplication, or reordering limits apply at the
Mission application? How much buffering is required, and what should
happen at saturation? Name the load fixture, observation point, and
reference stream for \ValidationRef{VAL-117}. Keep the 10 ms local
publication allocation distinct from mission delivery limits.

\paragraph{Q-121. Field platform coverage.}\label{Question:Q-121}
Which part of the subsystem must execute on macOS, iOS, Windows, and Linux?
Which operating-system versions, processor architectures, connection
methods, permissions, and drivers qualify on each? Which platforms are
required in the next baseline and which will be delivered later? Identify
the representative field devices and the connection and reading-receipt
success criteria for each platform. Reconcile the result with
\hyperref[Question:Q-106]{Q-106} and the macOS-only
\RequirementRef{REQ-C-103}.

\paragraph{Q-122. Readiness, validity, and stale state.}\label{Question:Q-122}
What edge or level constitutes one ready event, and how are multiple or
partial events handled? Which inputs classify a sample as valid, invalid,
or unqualified? Does an invalid or unqualified sample reset the stale
timer? What is reported before the first valid reading and at the 100 ms
boundary, including simultaneous arrival of a valid sample? What event
restores valid output after staleness or a fault? Provide expected
observations for the readiness and stale-state tests.

\subsection{Additional Decisions}

\paragraph{Q-123. Timing and measurement uncertainty.}\label{Question:Q-123}
Where are conversion-ready, timestamp capture, local publication, and
stale-state transition observed? How are corresponding events and records
matched? What clock synchronization, resolution, calibration, and maximum
measurement uncertainty are required? How should a result whose uncertainty
straddles the 10 ms, 100 ms, or 1 ms limit be classified? Identify the
approved test loads, fixtures, and repeat conditions.

\paragraph{Q-124. Endpoint replacement.}\label{Question:Q-124}
Which alternate sensor and Mission application qualify as compatible
replacements? Which configurations, adapters, or setup changes are permitted,
and which would count as redesign? What sample and status workflows must
succeed before and after each independent swap, and what record of the
changes establishes that the connection was reused?

\paragraph{Q-125. OMS status meaning.}\label{Question:Q-125}
Which combinations of connection and operating-health evidence map to valid,
stale, and faulted status? How are contradictory, ambiguous, or missing
observations handled, and which state has precedence when several apply?
For each approved state or transition, what should the gateway emit and what
should the Mission application infer? Record the applicable program values
and approval source in \hyperref[Table:OMSStatusDecisions]{the OMS decision
register} before advancing \RequirementRef{REQ-OMS-01}.

\paragraph{Q-126. Mission delivery degradation and recovery.}\label{Question:Q-126}
What should the Mission application observe during loss of sample input,
loss of OMS delivery, and restoration? During an outage, should the gateway
discard, retain, or coalesce samples and status, and how are gaps signaled?
What sample or status may be published first after restoration? Specify
loss, recovery-time, and freshness limits, and identify the test fixtures
and the observation points for both mission outputs.

\section{Recording and Playback Decisions}
\label{Section:RecordingPlaybackDecisions}
Owner: Michael Stoddard for each question below. Status: Open. Decisions require stakeholder review and the applicable acceptance-authority disposition before dependent implementation or quantitative verdicts.

\paragraph{Q-127. Recording boundary and fault policy.}\label{Question:Q-127}
Select the recording input boundary and schema, including received bytes or decoded records, invalid/partial inputs, sample/result associations, event order, start/stop boundaries, stop-and-flush behavior, file identity, capture-start processing state and any prehistory or checkpoint needed to reproduce it, known-gap representation, interrupted-file recovery, and storage-failure indications. Define retention, storage permissions, sensitive-data handling, and resource budgets; replay cannot recover physical or transfer observations that were not captured.

\paragraph{Q-128. Recording context and compatibility.}\label{Question:Q-128}
Define format/version identification, required software/configuration/calibration metadata, available versus unknown fields, context presentation, accepted compatibility combinations, treatment of missing initial-state information, corrupt/truncated-file classification, rejection reasons, and revalidation when a selected file changes. Identify compatible and incompatible fixtures and the expected disposition of each.

\paragraph{Q-129. Playback time, state, and controls.}\label{Question:Q-129}
Define the playback processing boundary, event tie-break and rollover rules, supported speeds relative to real time, the recorded-interval/speed timing reference and its treatment of speed-change and pause boundaries, timing-error limits in milliseconds, observation clock and uncertainty, pause/resume boundaries, valid seek locations, restoration strategy, initialization on every start/restart independent of prior live or replay sessions, transient-output disposition, single-step event meaning, presentation, and completion/failure behavior. Choose replay-clock semantics for freshness and timeouts during pause, stepping, speed changes, and seeking. Decide concurrency with live input; it is not assumed. Set state/output reference oracles, including recordings begun during operation and the disposition when their initial state cannot be established. Reset, captured-state restoration and prehistory replay remain alternatives pending approval. Software replay reproduces captured inputs, not the physical sensor or USB hardware.

\paragraph{Q-130. Replay identification and routing.}\label{Question:Q-130}
Agree how playback is identified at presentation and sample/status output boundaries, the test-consumer contract, endpoint selection and changes, handling of absent selection or delivery failure, and conditions for any mission endpoint. Reconcile with the OMS decisions; no unapproved OMS field or live mission delivery is selected by this proposal.

\paragraph{Q-131. Recording and playback acceptance budgets.}\label{Question:Q-131}
Approve sustained recording rate in records/s, burst rate and duration, minimum run duration in seconds or hours, storage allocation, recorder loss count/fraction and observation interval, allowed added live latency in milliseconds, allowed delivery-rate/gap degradation, and per-control response limits in milliseconds or seconds. Identify fixture sizes, concurrent loads, storage conditions, repeat counts, calibrated clocks and uncertainty/verdict rules. Agree the representative fault, stakeholder tasks and success criteria for the four use-case validations. These TBDs block quantitative verdicts; no numerical value is inferred from existing live-path timing allocations.

\section{Risks and Lifecycle Gates}
Unresolved capture coherence may pair a conversion with the wrong counter value. Undefined qualification or publication semantics may produce an unusable decision or a misinterpreted result. Source connection may be confused with source health. ADC0804 and FT245R are selected, but remaining electrical, timing, allocation, and macOS-version decisions still block detailed design and conformance claims. Owners and affected decisions are identified in \hyperref[Table:OwnedQuestions]{the question register}. No probabilities, severity ratings, schedules, or mitigations have been approved.

The documentation gate requires supported metadata, explicit gaps, consistent traceability, a complete build, visual inspection, and human review. Later needs, requirements, architecture, implementation, verification, validation, and software-release gates require their own evidence and authorized dispositions. Milestones remain \hyperref[Question:Q-110]{Q-110}.
